\begin{textsample}{POS Dim 3 – ba – Score 9.00 – ba/ba\_inf\_4.txt}  \label{ex:f3_pos_207}
7.4.
Articulações e Garantias É preciso afirmar a infância como “ um mundo outro ” ( CHOMBART DE LAUWE, 1991 ).
É preciso, portanto, não fazer da conjugação da Educação \textbf{Infantil} com o Ensino Fundamental um processo de subordinação a este e a sua tradição pedagógica, para não cairmos numa certa pedagogização essencialista da Educação \textbf{Infantil}.
Essa atitude que defendemos não significa, como alerta Kuhlmann Jr., uma demarcação no estilo avestruz, face à realidade nacional, fazendo uma reserva de domínio descomprometida com o Ensino Fundamental.
Trata -se de reconhecer a importância dos dois momentos, que deverão entretecer -se, articular -se intensamente, pois vivenciar a Educação \textbf{Infantil} e suas especificidades e valorizar o Ensino Fundamental como uma demanda de potência democrática ineliminável é reconhecer que, afinal de contas, formação, aqui, está tratando da infância na sua condição humana e nas suas diferenças ontológicas.
Ao publicizar o Programa de Ampliação do Ensino Fundamental, o Ministério da Educação aponta para uma ação de equidade social, ou seja, de justiça social pela via da diferença, ao garantir a todas as \textbf{crianças} mais um ano de escolarização, corrigindo, por consequência, as distorções produzidas pelos mecanismos do Fundef.
Não se tratava tão somente da antecipação da escolarização obrigatória, mas também de sua ampliação, o que poderia significar mais permanência com mais qualidade.
No que se refere à Educação \textbf{Infantil}, apostou -se que o deslocamento das \textbf{crianças} de 6 anos para o Ensino Fundamental geraria, como efeito positivo, o aumento de vagas nessa etapa educativa.
Em realidade, estudos de Campos e Silva ( 2011 ) mostram que a transição tem sido feita de forma inconsistente em alguns contextos educacionais brasileiros, e que as expectativas da iniciativa não vêm ocorrendo.
Por exemplo, a pressão por vagas na Educação \textbf{Infantil} continua.
Preocupar -se com uma prática pedagógica que não compartimentalize Educação \textbf{Infantil} e Ensino Fundamental é fundar mais uma atitude de desconstrução das fragmentações que tanto aprendemos a fabricar.
Aproximar, articulando, esses dois momentos educacionais é um procedimento que nos direciona justamente para a complexidade da aprendizagem como fenômeno sociointerativo que deve sempre, de forma inarredável, estar compromissada com a qualificação da formação.
É preciso que qualquer orientação norteadora sobre a educação das \textbf{crianças} \textbf{pequenas} parta das suas ontologias, das suas condições socioculturais concretas, compreendendo que o conhecimento do mundo envolve \textbf{afeto}, prazer, desprazer, fantasia, \textbf{brincadeira}, movimento, poesia, ciências, artes, linguagem, música, matemática etc., que, para a \textbf{criança}, a \textbf{brincadeira} é uma forma de linguagem, assim como a linguagem pode ser uma forma de \textbf{brincadeira} ( KUHLMANN JR., 1999 ), e que o polimorfismo é a forma de expressão preponderante deste mundo outro.
Nessa nova realidade, o Ensino Fundamental de 9 ( nove ) anos traz diversas implicações, desde a exigência de investimentos em aspectos físicos e materiais para atender à \textbf{criança} de 6 ( seis ) anos no espaço escolar, até a destinação de recursos relacionados às políticas de formação de professores e demais necessidades de cunho pedagógico.
Podemos dizer que, mesmo com o avanço das políticas públicas para a \textbf{pequena} infância, verificamos que as propostas educativas criadas para este momento histórico revelam que, de forma elaborada, o caráter compensatório ainda se faz presente.
Vale ressaltar que a Educação \textbf{Infantil} é complementar à educação da família e um direito da \textbf{criança} ; portanto não se constitui apenas como um pré-requisito para o seu ingresso na escola.
Reiteramos que o fato de a Educação \textbf{Infantil} fazer parte da Educação Básica não subentende que deverá antecipar o modelo escolar do Ensino Fundamental, extremamente contestável quando se trata da relação dessa proposta pedagógica com as diversas infâncias que se apresentam à escola buscando formação.
Convencido da necessidade de parâmetros para a educação da infância, o MEC ( BRASIL, 1995, p. 11 ) apresentou critérios de qualidade para que a Educação \textbf{Infantil} respeitasse os direitos fundamentais das \textbf{crianças} : - Nossas \textbf{crianças} têm direito à \textbf{brincadeira} ; - Nossas \textbf{crianças} têm direito à atenção individual ; - Nossas \textbf{crianças} têm direito a um ambiente aconchegante, seguro e estimulante ; - Nossas \textbf{crianças} têm direito ao contato com a natureza ; - Nossas \textbf{crianças} têm direito à \textbf{higiene} e à saúde ; - Nossas \textbf{crianças} têm direito a uma alimentação sadia ; - Nossas \textbf{crianças} têm direito a desenvolver sua \textbf{curiosidade}, imaginação e capacidade de expressão ; - Nossas \textbf{crianças} têm direito ao movimento em espaços amplos ; - Nossas \textbf{crianças} têm direito à proteção, ao \textbf{afeto} e à amizade ; - Nossas \textbf{crianças} têm direito a expressar seus sentimentos ; - Nossas \textbf{crianças} têm direito a uma especial atenção durante seu período de adaptação à \textbf{creche} ; - Nossas \textbf{crianças} têm direito a desenvolver sua identidade cultural, racial e religiosa.
O que podemos concluir dessas elaborações de possibilidades autonomistas e emancipacionistas é que devemos falar não apenas de direito à aprendizagem, mas de direito à formação qualificada via aprendizagens relevantes e pertinentes.
Toda proposta curricular que tem como base propostas de aprendizados precisa ser submetida a processos ampliados de debate sobre sua valoração, pois esse é o campo do formacional.
Na medida em que nem toda aprendizagem é boa, a formação requer uma aprendizagem eleita como formativa, do contrário podemos estar repetindo padrões e crenças legitimados por centros de poder da educação, sem reflexões profundas sobre suas pertinências e relevâncias formativas, até porque, levando em conta os argumentos aqui já desenvolvidos sobre a formação, a formação, da perspectiva curricular, é sempre uma (in)tensa pauta ética e política, e envolve, portanto, opções, escolhas em face de demandas diversas.
Nesse mesmo veio, devemos nos interrogar quem elege o que é formativo para as \textbf{crianças}, como elegem e para que elegem.
Vale dizer que, para nosso interesse de texto, o formativo para as \textbf{crianças} implica compreender o que e como elas estão aprendendo, envolvendo os atos de currículo ( MACEDO, 2010 ; 2012 ; 2013 ) que envolvem e fazem a mediação dessa aprendizagem.
Portanto, não falemos apenas de direito à aprendizagem, mas de direito à formação qualificada, em que aprendizagens mediadas e com conteúdos cognitivos, éticos, estéticos, políticos e culturais estejam presentes em processos de qualificação dessa própria formação.
Vale ressaltar o argumento de Kohan ( 2003 ).
Para esse autor, em educação, associamos infância à primeira idade, pensando os seres humanos atravessando estágios cujos percursos costumam ter o signo do progresso.
A infância seria o primeiro degrau associado a uma marca do ser em potência – e não em ato –, do que pode ser, mas ainda não é, do que virá a ser se acompanhamos a infância com um bom currículo, numa referência aberta a uma concepção exterodeterminante de desenvolvimento, a uma perspectiva de quase condenação sociopedagógica.
A proposta pedagógica para a Educação \textbf{Infantil} e a incorporação da \textbf{criança} de 6 anos no Ensino Fundamental é justificada como uma tentativa de superar um discurso que difundiu o assistencialismo separado da educação.
Buscou -se, e ainda se busca, a construção de uma identidade localizada em outro extremo, exaltando a excelência educativa pautada no modelo escolar, como se este fosse o melhor parâmetro de qualidade.
A permanência de \textbf{crianças} de 6 anos no Ensino Fundamental hoje é fato, ou seja, é uma política que se efetivou.
Com efeito, em 2006, a redação da Lei no 11.274 modificou o art. 32 da Lei de Diretrizes e Bases da Educação Nacional de 1996.
A partir de então, definiu -se que o aluno ingressaria nos estabelecimentos escolares que ministravam o Ensino Fundamental aos 6 anos de idade, em caráter obrigatório.
Alguns incisos do artigo mencionado afirmam a necessidade de que os estudantes desenvolvam, ao longo da segunda etapa da Educação Básica, a habilidade de aprender, considerando relevante o pleno domínio da leitura, da escrita e do cálculo, além da percepção adequada do ambiente natural e social.
O estudo do sistema político, da tecnologia, das artes e dos valores que regem a sociedade em que vivemos deveria fazer parte do currículo.
O Ministério da Educação ( MEC ), a Secretaria de Educação Básica ( SEB ), o Departamento de Políticas de Educação \textbf{Infantil} e Ensino Fundamental ( DPE ) e a Coordenação Geral do Ensino Fundamental ( COEF ) publicaram, em 2004, um documento intitulado “ Ensino Fundamental de nove anos : orientações gerais ”, dizendo ter como objetivo incentivar políticas que pudessem promover transformações estruturais nas instituições escolares, no que se refere ao processo de “ ensino/aprendizagem, avaliação, currículo, conhecimento e desenvolvimento humano ”.
Ao mesmo tempo, postulava -se a necessidade de que não houvesse uma ruptura entre a Educação \textbf{Infantil} e o Ensino Fundamental e, nessa direção, reforçava -se a importância das Diretrizes Curriculares Nacionais para a Educação \textbf{Infantil} na revisão da proposta pedagógica do Ensino Fundamental, que passaria a atender às \textbf{crianças} de 6 anos : > “ [... ] O objetivo de um maior número de anos de ensino obrigatório é assegurar a todas as \textbf{crianças} um tempo mais longo de convívio escolar, maiores oportunidades de aprender e, com isso, uma aprendizagem mais ampla.
É evidente que a maior aprendizagem não depende do tempo de permanência na escola, mas sim do emprego mais eficaz do tempo.
No entanto, a associação de ambos deve contribuir, significativamente, para que os estudantes aprendam mais.
Seu ingresso no Ensino Fundamental obrigatório não pode constituir -se em medida meramente administrativa.
O \textbf{cuidado} na sequência do processo de desenvolvimento e aprendizagem das \textbf{crianças} de seis anos de idade implica o conhecimento e a atenção às suas características etárias, sociais e psicológicas.
As orientações pedagógicas, por sua vez, estarão atentas a essas características para que as \textbf{crianças} sejam respeitadas como sujeitos do aprendizado ” ( BRASIL, 2004, p. 17-18 ).
Importante notar que, se no âmbito formal o primeiro ano passou a pertencer ao Ensino Fundamental, do ponto de vista da prática pedagógica ainda está inserido na Educação \textbf{Infantil}.
Por isso, as Diretrizes continuam sendo orientações para o planejamento dos conteúdos que deveriam ser trabalhados com as \textbf{crianças} de 6 anos.
Recomenda -se atenção às singularidades dos alunos dessa faixa etária, a não antecipação do currículo da antiga primeira série, ao mesmo tempo em que se estimula a alfabetização precoce.
O fato de o educando estar imerso em um “ ambiente alfabetizador ” já na Educação \textbf{Infantil}, ou seja, de ter acesso a situações em que a leitura e a escrita possuem usos reais de expressão e comunicação, seria um elemento facilitador para um processo de transição “ natural ” entre a primeira e a segunda etapas da Educação Básica.

% matched lemmas: afeto, brincadeira, creche, criança, cuidado, curiosidade, higiene, infantil, pequeno
\end{textsample}
